% Generated by scripts/assemble.py; corrections belong in source/corrections.json.
\examhead{131}{武汉大学 2025 年数学分析试卷}{124}
% p125-b22; source PDF page 125
\par\Needspace{5\baselineskip}
\blockanchor{p125-b22}
\majorquestion{一}计算题.\par
% p125-b23; source PDF page 125
\blockanchor{p125-b23}
\question{1}求极限  $\lim_{n\to\infty}\cos\frac{\pi}{2^{2}}\cos\frac{\pi}{2^{3}}\cdots\cos\frac{\pi}{2^{n}}$ .\par
% p125-b24; source PDF page 125
\blockanchor{p125-b24}
\question{2}已知 $\sin (xy) + 2y^2 = 1$，求 $\frac{\mathrm{dy}}{\mathrm{dx}}$.\par
% p125-b25; source PDF page 125
\blockanchor{p125-b25}
\question{3}$I_{n} = \int_{0}^{\frac{\pi}{4}}(\sin x)^{n}\cos x\mathrm{d}x,$ 求 $\sum_{n = 0}^{\infty}I_n.$\par
% p125-b26; source PDF page 125
\par\Needspace{5\baselineskip}
\blockanchor{p125-b26}
\majorquestion{二}计算题.\par
% p125-b27; source PDF page 125
\blockanchor{p125-b27}
\question{1}已知 $D$ 是由 $x = 0, y = \frac{\pi}{2}, x = y^2$ 所围的区域, 求 $\iint_{D} \frac{\sin y}{y} \, \mathrm{d}x \, \mathrm{d}y$.\par
% p126-b00; source PDF page 126
\par\Needspace{6\baselineskip}
\blockanchor{p126-b00}
\question{2}$\Sigma$ 是 $y = \sqrt{x} (0 \leq x \leq 1)$ 绕 $x$ 轴旋转形成的曲面, 法向量方向与 $x$ 轴正方向成钝角, 求第二型曲面积分\par
\nopagebreak[4]
% p126-b01; source PDF page 126
\blockanchor{p126-b01}
\displayformula{I = \iint_ {\Sigma} 2 (x + y) \mathrm{d} y \mathrm{d} z + y z \mathrm{d} x \mathrm{d} y.}
% p126-b02; source PDF page 126
\blockanchor{p126-b02}
\question{3}已知 $\varphi(t) = \int_0^{+\infty} \frac{\ln(1 + tx)}{x(1 + x)} \, \mathrm{d}x$, 求 $\varphi(1)$.\par
% p126-b03; source PDF page 126
\par\Needspace{5\baselineskip}
\blockanchor{p126-b03}
\majorquestion{三}证明题.\par
% p126-b04; source PDF page 126
\blockanchor{p126-b04}
\question{1}已知 $p, q > 0$, 且 $\frac{1}{p} + \frac{1}{q} = 1$, 求证: 对任意正数 $x_{1}, x_{2}$, 有 $x_{1}x_{2} \leq \frac{1}{p} x_{1}^{p} + \frac{1}{q} x_{2}^{q}$.\par
% p126-b05; source PDF page 126
\par\Needspace{6\baselineskip}
\blockanchor{p126-b05}
\question{2}$f(x)$ 在 $(- \infty, + \infty)$ 上二阶可微，且 $f(x) > 0$ ，$M = \sup_{x \in \mathbb{R}} |f''(x)| < + \infty$。证明：\par
\nopagebreak[4]
% p126-b06; source PDF page 126
\blockanchor{p126-b06}
\displayformula{(f ^ {\prime} (x)) ^ {2} \leq 2 M f (x), x \in (- \infty , + \infty).}
% p126-b07; source PDF page 126
\blockanchor{p126-b07}
\question{3}已知 $f(x)$ 在 $[a, b]$ 上连续, 并且满足对于 $i = 0, 1, 2$, 有 $\int_{a}^{b} x^{i} f(x) \mathrm{d}x = 0$. 求证: 存在三个不同的数 $x_{1}, x_{2}, x_{3} \in (a, b)$, 使得 $f(x_{1}) = f(x_{2}) = f(x_{3}) = 0$.\par
% p126-b08; source PDF page 126
\blockanchor{p126-b08}
\question{4}$\alpha$ 是正数, 判断函数项级数 $\sum_{n=0}^{\infty} x^{\alpha} e^{-nx}$ 在 $(0, +\infty)$ 上是否一致收敛.\par
% p126-b09; source PDF page 126
\blockanchor{p126-b09}
\question{5}已知 $f(x, y)$ 在 $x^2 + y^2 < 1$ 上三阶连续可导，并且 $f(0, 0) = 0$，求证：存在二阶连续可微函数 $g_1(x, y), g_2(x, y)$，使得 $f(x, y) = xg_1(x, y) + yg_2(x, y)$.\par
% p126-b10; source PDF page 126
\blockanchor{p126-b10}
\question{6}求 $(a, b)$ 的两个区域, 使得 $f_{a, b}(x, y) = ay^2 + bx$ 在约束条件 $x^2 + y^2 = 1$ 下分别有两个和四个临界点.\par
